% Part: incompleteness
% Chapter: theories-computability
% Section: first-incompleteness

\documentclass[../../../include/open-logic-section]{subfiles}

\begin{document}

\olfileid{inc}{tcp}{inc}

\olsection{$\Th{Q}$ 没有完备、一致、可公理化的扩张}

\begin{thm}
\ollabel{thm:first-incompleteness}
不存在 $\Th{Q}$ 的完备、一致且可公理化的扩张。
\end{thm}

\begin{proof}
我们已经知道，不存在 $\Th{Q}$ 的一致且可判定的扩张。但如果 $\Th{T}$ 是完备且已公理化的，那么它是可判定的。 
\end{proof}

\begin{explain}
这个定理与 哥德尔 在 1931 年对第一不完全性定理的原始表述相去不远。除了更现代的术语，关键区别在于：哥德尔 使用的是“$\omega$-一致的”而非“一致的”；并且他无法完全一般性地谈论“可公理化”，因为当时可计算性的形式概念尚未建立。（可计算性的形式模型是在接下来的十年中发展起来的，其中包括 哥德尔 的工作，很大程度上正是为了刻画易受 哥德尔 现象影响的理论。）

这个定理说的是你无法三者兼得，即完备性、一致性和可公理性。不过，如果你放弃其中任何一个，就可以得到另外两个：$\Th{Q}$ 是一致的且可计算公理化，但不完备；不一致的理论是完备的且可计算公理化（例如，以 $\{ \eq/[0][0] \}$ 为公理），但不一致；而算术真语句的集合是完备且一致的，但不可计算公理化。
\end{explain}

\end{document}